\documentclass{beamer}

% ----- PACKAGES -------
\usepackage{transparent}
\usepackage{xcolor}
\usepackage{multicol}
\usepackage{listings}
\usepackage{hyperref}

%-----------------------


\title[Introducción a LaTeX]{\LaTeX: Cómo escribir TeXtos bonitos con \textbackslash emph\{comandos feos\}}

\subtitle{Introducción a LaTeX \\ (eso que necesitas para el TFG que aún no sabes) :3}

\author[Feito, Herrador \& Lúa]{Daniel Feito, Hugo Herrador \& Lúa Rico}
\institute{DAFIC $\times$ GPUL}
\date[GPUL School 2025]{Marzo de 2025}

\usetheme{Madrid}
\usecolortheme[RGB={50,110,118}]{structure}
\setbeamertemplate{navigation symbols}{}


\begin{document}

\frame{\titlepage}

\AtBeginSection[]
{
  \begin{frame}
    \frametitle{Índice}
    \begin{multicols}{2}
    \tableofcontents[currentsection]
    \end{multicols}
  \end{frame}
}

\begin{frame}[fragile]
\frametitle{Repositorio de las diapositivas}

\begin{figure}
    \centering
    \includegraphics[width=0.5\linewidth]{images/QR.png}
\end{figure}

\end{frame}


\input{content/introducción}
\section{Tipos de documentos}

\begin{frame}[fragile]
\frametitle{Tipos de documentos}

\begin{itemize}
 \item \verb|\document{article}| $\rightarrow$ Artículos científicos, informes breves, documentación de programas
 \item \verb|\document{book}| $\rightarrow$ Libros completos con capítulos
 \item \verb|\document{report}| $\rightarrow$ Tesis / informes técnicos extensos
 \item \verb|\document{letter}| $\rightarrow$ Cartas
 \item \verb|\document{beamer}| $\rightarrow$ Diapositivas 

\end{itemize}

\end{frame}

% Comandos básicos

% Crear portada (Campos básicos)

\section{Comandos básicos}

\begin{frame}[fragile]
\frametitle{Comandos básicos: portada}

Opción fácil $\rightarrow$ \verb|\maketitle|
\begin{itemize}
     \item \verb|\title{Titulo documento}|
     \item \verb|\author{Daniel Feito \& Hugo Herrador \& Lúa Rico}|
     \item \verb|\date{March 2025}|
\end{itemize}

\vspace{0.5cm}

Opción avanzada $\rightarrow$ \verb|\begin{titlepage}|


\end{frame}

% Section, subsection
\begin{frame}[fragile]
\frametitle{Comandos básicos: secciones, índice y páginas}
Crear secciones y subsecciones
\begin{itemize}
    \item \verb|\section{}|
    \item \verb|\subsection{}|
    \item \verb|\subsubsection{}| 
\end{itemize}

\vspace{0.5cm}

Crear el índice (se actualiza automáticamente) $\rightarrow$ \verb|\tableofcontents|

\vspace{0.5cm}

Crear una nueva página  $\rightarrow$ \verb|\newpage|

\end{frame}

% Cargar paquetes

\begin{frame}[fragile]
\frametitle{Comandos básicos: cargar paquetes}
Normalmente se hace justo después de definir el tipo de documento. 

\vspace{0.5cm}

\verb|\usepackage{nombre_paquete}|
\end{frame}

% Negrita, cursiva, colores en texto
\begin{frame}[fragile]
\frametitle{Comandos básicos: negrita, cursiva y colores}
\begin{itemize}
    \item \textbf{Negrita} $\rightarrow$ \verb|\textbf{Negrita}|
    \item \textit{Cursiva} $\rightarrow$ \verb|\textit{Cursiva}|
    \begin{itemize}
        \item [*] Otra forma de escribir cursiva es con el comando \verb|\emph{}| (`enfatizar' texto)
        \item [*] La diferencia es que si se usa 2 veces \verb|\emph{}|, el texto vuelve a escribirse normal, pero si se usa 2 veces \verb|\textit{}| el texto se queda en cursiva
    \end{itemize}
    \item \textcolor{red}{Colores} $\rightarrow$ \verb|\textcolor{red}{Colores}| (necesario el paquete \verb|\usepackage{xcolor}|)
\end{itemize}

\end{frame}


\section{Listas}

\begin{frame}[fragile]
\frametitle{Listas: Tipos y formato}

Tipos de listas:
\begin{itemize}
 \item Listas ordenadas $\rightarrow$ \verb|enumerate|
 \item Listas no ordenadas $\rightarrow$ \verb|itemize|
 \item Listas de descripciones $\rightarrow$ \verb|description|
\end{itemize}

\vspace{0.5cm}

Cada elemento se separa con \verb|\item| .

\vspace{0.5cm}

Las listas pueden anidarse (aparecer una dentro de otra), incluso entre tipos distintos.
\end{frame}


\begin{frame}[fragile]
\frametitle{Listas: Ejemplo de lista ordenada}
\begin{columns}
\column{0.45\textwidth}

    \begin{enumerate}
    \item Item A
    \item Item B
        \begin{enumerate}
            \item Item B.1
            \item Item B.2
        \end{enumerate}
    \item Item C
    \item Item D
    \end{enumerate}

\column{0.45\textwidth}
\begin{semiverbatim}
\\begin\{enumerate\}
    \\item Item A
    \\item Item B
        \\begin\{enumerate\}
            \\item Item B.1
            \\item Item B.2
        \\end\{enumerate\}
    \\item Item C
    \\item Item D
\\end\{enumerate\}
\end{semiverbatim}
\end{columns}
\end{frame}


\begin{frame}[fragile]
\frametitle{Listas: Ejemplo de lista no ordenada}
\begin{columns}
\column{0.45\textwidth}

    \begin{itemize}
    \item Item A
    \item Item B
        \begin{itemize}
            \item Item B.1
            \item Item B.2
        \end{itemize}
    \item Item C
    \item Item D
    \end{itemize}

\column{0.45\textwidth}
\begin{semiverbatim}
\\begin\{itemize\}
    \\item Item A
    \\item Item B
        \\begin\{itemize\}
            \\item Item B.1
            \\item Item B.2
        \\end\{itemize\}
    \\item Item C
    \\item Item D
\\end\{itemize\}
\end{semiverbatim}
\end{columns}
\end{frame}

\begin{frame}[fragile]
\frametitle{Listas: Ejemplo de lista de descripciones}
\begin{description}
   \item[Término A] Descripción súper ingeniosa aquí.
   \item[Término B] Otra descripción creativa aquí.
\end{description}

\vspace{0.5cm}

\begin{semiverbatim}
\\begin\{description\}
   \\item[Término A] Descripción súper ingeniosa aquí.
   \\item[Término B] Otra descripción creativa aquí.
\\end\{description\}
\end{semiverbatim}
\end{frame}

\begin{frame}[fragile]
\frametitle{Listas: Estilos}
Estilos de las listas:
\begin{itemize}
    \item Las listas ordenadas por defecto aparecen con números.
    \item Las listas no ordenadas por defecto aparecen con puntos.
    \item El estilo de las listas anidadas puede cambiar en cada nivel de profundidad.
    \item El decorador de cada \verb|\item| puede personalizarse indicándolo entre corchetes:
\end{itemize}

\begin{columns}
\column{0.4\textwidth}

\begin{itemize}
    \item[-] Ejemplo con `-'.
    \item[!] Ejemplo con `!'.
    \item[XD] Ejemplo con texto.
    \item[] Ejemplo sin nada.
\end{itemize}

\column{0.65\textwidth}
\begin{semiverbatim}
\\begin\{itemize\}
    \\item[-] Ejemplo con `-'.
    \\item[!] Ejemplo con `!'.
    \\item[XD] Ejemplo con texto.
    \\item[] Ejemplo sin nada.
\\end\{itemize\}
\end{semiverbatim}
\end{columns}

\end{frame}

\input{content/imágenes}
\section{Enlaces}
\section{Tablas}

\begin{frame}
\frametitle{Tablas}

\centering
\url{https://www.tablesgenerator.com/}
    
\end{frame}
\input{content/código}
\section{Matemáticas}

\begin{frame}[fragile]
    Calculemos la siguiente integral definida:
    \begin{equation}
        \int_{- \infty}^{+ \infty}{e^{-x^2}dx}
    \end{equation}
    \begin{verbatim}
    \begin{equation}
        \int_{- \infty}^{+ \infty}{e^{-x^2}dx}
    \end{equation}
    \end{verbatim}
    \transparent{0.3}
    Otra opción: $\int_{- \infty}^{+ \infty}{e^{-x^2}dx}$
    \begin{verbatim}
        $\int_{- \infty}{+ \infty}{e^{-x^2}dx}$
    \end{verbatim}
    Tercera opción: $$\int_{- \infty}^{+ \infty}{e^{-x^2}dx}$$
    \begin{verbatim}
        $$\int_{- \infty}{+ \infty}{e^{-x^2}dx}$$
    \end{verbatim}
\end{frame}

\begin{frame}[fragile]
    Calculemos la siguiente integral definida:
    \transparent{0.3}
    \begin{equation*}
        \int_{- \infty}^{+ \infty}{e^{-x^2}dx}
    \end{equation*}
    \begin{verbatim}
    \begin{equation}
        \int_{- \infty}^{+ \infty}{e^{-x^2}dx}
    \end{equation}
    \end{verbatim}
    \transparent{1}
    Otra opción: $\int_{- \infty}^{+ \infty}{e^{-x^2}dx}$
    \begin{verbatim}
        $\int_{- \infty}{+ \infty}{e^{-x^2}dx}$
    \end{verbatim}
    \transparent{0.3}
    Tercera opción: $$\int_{- \infty}^{+ \infty}{e^{-x^2}dx}$$
    \begin{verbatim}
        $$\int_{- \infty}{+ \infty}{e^{-x^2}dx}$$
    \end{verbatim}
\end{frame}

\begin{frame}[fragile]
    Calculemos la siguiente integral definida:
    \transparent{0.3}
    \begin{equation*}
        \int_{- \infty}^{+ \infty}{e^{-x^2}dx}
    \end{equation*}
    \begin{verbatim}
    \begin{equation}
        \int_{- \infty}^{+ \infty}{e^{-x^2}dx}
    \end{equation}
    \end{verbatim}
    \transparent{0.3}
    Otra opción: $\int_{- \infty}^{+ \infty}{e^{-x^2}dx}$
    \begin{verbatim}
        $\int_{- \infty}{+ \infty}{e^{-x^2}dx}$
    \end{verbatim}
    \transparent{1}
    Tercera opción: $$\int_{- \infty}^{+ \infty}{e^{-x^2}dx}$$
    \begin{verbatim}
        $$\int_{- \infty}{+ \infty}{e^{-x^2}dx}$$
    \end{verbatim}
\end{frame}

\begin{frame}[fragile]
    \begin{center}
    Sea $\mathcal{I} = \int_{- \infty}^{+ \infty}{e^{-x^2}dx}$
    \end{center}
    \transparent{0}
    \begin{align*}
        \mathcal{I}^2 &= \left ( \int_{- \infty}^{+ \infty}{e^{-x^2}dx} \right )^2 \\
        &= \left ( \int_{- \infty}^{+ \infty}{e^{-x^2}dx} \right ) \cdot \left ( \int_{- \infty}^{+ \infty}{e^{-y^2}dy} \right ) \\
        &= \left ( \int_{- \infty}^{+ \infty}{e^{-(x^2 + y^2)}dx} \right )
    \end{align*}
    \begin{center}
    ¿Qué podemos hacer ahora?
    \end{center}
    \begin{center}
    ¡Un cambio a polares!
    \end{center}
    $$
    \begin{cases}
        x = r \cdot \cos{\theta} \\
        y = r \cdot \sin{\theta}
    \end{cases}
    $$
\end{frame}

\begin{frame}[fragile]
    \begin{center}
    Sea $\mathcal{I} = \int_{- \infty}^{+ \infty}{e^{-x^2}dx}$
    \end{center}
    \begin{align*}
        \mathcal{I}^2 &= \left ( \int_{- \infty}^{+ \infty}{e^{-x^2}dx} \right )^2 \\
        &= \left ( \int_{- \infty}^{+ \infty}{e^{-x^2}dx} \right ) \cdot \left ( \int_{- \infty}^{+ \infty}{e^{-y^2}dy} \right ) \\
        &= \left ( \int_{- \infty}^{+ \infty}{e^{-(x^2 + y^2)}dx} \right )
    \end{align*}
    \transparent{0}
    \begin{center}
    ¿Qué podemos hacer ahora?
    \end{center}

    \begin{center}
    ¡Transformar a polares!
    \end{center}
    $$
    \begin{cases}
        x = r \cdot \cos{\theta} \\
        y = r \cdot \sin{\theta}
    \end{cases}
    $$
\end{frame}

\begin{frame}[fragile]
    \begin{center}
    Sea $\mathcal{I} = \int_{- \infty}^{+ \infty}{e^{-x^2}dx}$
    \end{center}
    \begin{align*}
        \mathcal{I}^2 &= \left ( \int_{- \infty}^{+ \infty}{e^{-x^2}dx} \right )^2 \\
        &= \left ( \int_{- \infty}^{+ \infty}{e^{-x^2}dx} \right ) \cdot \left ( \int_{- \infty}^{+ \infty}{e^{-y^2}dy} \right ) \\
        &= \left ( \int_{- \infty}^{+ \infty}{e^{-(x^2 + y^2)}dx} \right )
    \end{align*}
    \begin{center}
    ¿Qué podemos hacer ahora?
    \end{center}
    \transparent{0}
    \begin{center}
    ¡Transformar a polares!
    \end{center}
    $$
    \begin{cases}
        x = r \cdot \cos{\theta} \\
        y = r \cdot \sin{\theta}
    \end{cases}
    $$
\end{frame}

\begin{frame}[fragile]
    \begin{center}
    Sea $\mathcal{I} = \int_{- \infty}^{+ \infty}{e^{-x^2}dx}$
    \end{center}
    \begin{align*}
        \mathcal{I}^2 &= \left ( \int_{- \infty}^{+ \infty}{e^{-x^2}dx} \right )^2 \\
        &= \left ( \int_{- \infty}^{+ \infty}{e^{-x^2}dx} \right ) \cdot \left ( \int_{- \infty}^{+ \infty}{e^{-y^2}dy} \right ) \\
        &= \left ( \int_{- \infty}^{+ \infty}{e^{-(x^2 + y^2)}dx} \right )
    \end{align*}
    \begin{center}
    ¿Qué podemos hacer ahora?
    \end{center}

    \begin{center}
    ¡Transformar a polares!
    \end{center}
    $$
    \begin{cases}
        x = r \cdot \cos{\theta} \\
        y = r \cdot \sin{\theta}
    \end{cases}
    $$
\end{frame}

\begin{frame}[fragile]
    \begin{align*}
        \mathcal{I}^2 &= \int_{0}^{2 \pi}\int_0^{+ \infty} \begin{vmatrix} \frac{\partial x}{\partial r} && \frac{\partial x}{\partial \theta} \\ \frac{\partial y}{\partial r} && \frac{\partial y}{\partial \theta}\end{vmatrix} \cdot e^{-r^2} dr d\theta \\
        &= \int_{0}^{2 \pi}\int_0^{+ \infty} \begin{vmatrix} \cos \theta && -r \sin \theta \\ \sin \theta && r \cos \theta \end{vmatrix} \cdot e^{-r^2} dr d\theta \\
        &= \int_{0}^{2 \pi}\int_0^{+ \infty} r \cdot e^{-r^2} dr d\theta \\
        &= -\frac{1}{2} \cdot \int_{0}^{2 \pi}\int_0^{+ \infty} r \cdot -2r \cdot e^{-r^2} dr d\theta \\
        &= -\frac{1}{2} \cdot \int_{0}^{2 \pi} \left . e^{-r^2} \right |_{0}^{+ \infty} d\theta\\
        &= -\frac{1}{2} \cdot \int_{0}^{2 \pi} - d\theta \\
        &= \frac{1}{2} \cdot 2\pi \\
        &= \pi \implies \boxed{\mathcal{I} = \sqrt{\pi}}
    \end{align*}
\end{frame}

\begin{frame}[fragile]
\frametitle{Ajustar tamaño de delimitadores}
    \begin{center}
    \begin{equation*}
    \mathcal{I}^2 = \left ( \int_{- \infty}^{+ \infty}{e^{-x^2}dx} \right )^2
    \end{equation*}
    \begin{verbatim}
    \mathcal{I}^2 =
    \left ( \int_{- \infty}^{+ \infty}
    {e^{-x^2}dx} \right )^2
    \end{verbatim}
    \end{center}
    \begin{itemize}
        \item{\verb|\left| y \verb|\right| se usan para ajustar el tamaño de los delimitadores} 
        \item{Podemos usar . si sólo necesitamos uno de los dos lados}
        \item{Para los subíndices y superíndices usamos \_ y \^{}, respectivamente}
    \end{itemize}
\end{frame}


\begin{frame}[fragile]
\frametitle{Alinear fórmulas respecto del símbolo =}

\begin{align*}
    \mathcal{I}^2 &= \left ( \int_{- \infty}^{+ \infty}{e^{-x^2}dx} \right )^2 \\
    &= \left ( \int_{- \infty}^{+ \infty}{e^{-x^2}dx} \right ) \cdot \left ( \int_{- \infty}^{+ \infty}{e^{-y^2}dy} \right )
\end{align*}

\begin{verbatim}
\begin{align*}
    \mathcal{I}^2 &= \left ( \int {e^{-x^2}dx} \right )^2 \\
    &= \left ( \int {e^{-x^2}dx} \right ) \cdot
       \left ( \int{e^{-y^2}dy} \right )
\begin{align}
\end{verbatim}
\end{frame}

\begin{frame}[fragile]
    \frametitle{Símbolos varios}
        \begin{align*}
            \lim_{x\to x_{0}}f(x) = l :&\iff \Big[  \forall\epsilon>0, \exists\delta>0 \Big/ 0<\left|x-x_0\right|<\delta \\
            &\implies \left|f(x)-l\right| <\epsilon \Big]
        \end{align*}
        \begin{itemize}
            \item{Letras griegas (\textbackslash epsilon, \textbackslash delta...)}
            \item{Operadores lógicos (\textbackslash land, \textbackslash lor...)}
            \item{Implicaciones (\textbackslash implies, \textbackslash iff)}
            \item{Cuantificador existencial y universal (\textbackslash forall \textbackslash exists)}
            \item{Operadores de conjuntos (\textbackslash cup, \textbackslash cap, \textbackslash subset, \textbackslash in, \textbackslash notin...)}
        \end{itemize}
    \transparent{0}
    \center{\textbf{\textcolor{red}{RETO}: ¿NOS ATREVEMOS A COPIAR LA DEFINICIÓN DE LÍMITE?}}
\end{frame}

\begin{frame}[fragile]
    \frametitle{Símbolos varios}
        \begin{align*}
            \lim_{x\to x_{0}}f(x) = l :&\iff \Big[  \forall\epsilon>0, \exists\delta>0 \Big/ 0<\left|x-x_0\right|<\delta \\
            &\implies \left|f(x)-l\right| <\epsilon \Big]
        \end{align*}
        \begin{itemize}
            \item{Letras griegas (\textbackslash epsilon, \textbackslash delta...)}
            \item{Operadores lógicos (\textbackslash land, \textbackslash lor...)}
            \item{Implicaciones (\textbackslash implies, \textbackslash iff)}
            \item{Cuantificador existencial y universal (\textbackslash forall \textbackslash exists)}
            \item{Operadores de conjuntos (\textbackslash cup, \textbackslash cap, \textbackslash subset, \textbackslash in, \textbackslash notin...)}
        \end{itemize}
    \center{\textbf{\textcolor{red}{RETO}: ¿NOS ATREVEMOS A COPIAR LA DEFINICIÓN DE LÍMITE?}}
\end{frame}


\section{Columnas}

\begin{frame}[fragile]
\frametitle{Columnas}
\begin{itemize}
    \item Multicols $\rightarrow$ \verb|\begin{multicols}{número_columnas}| 
    \item Minipage  $\rightarrow$ \verb|\begin{minipage} ... \end{minipage}| 
\end{itemize}


% Minipage (imágenes): 
%     \begin{minipage}{0.45\textwidth}
%         \centering
%         \includegraphics[width=1\linewidth]{UDC-Logo.png} 
%     \end{minipage}\hfill
%     \begin{minipage}{0.45\textwidth}
%         \centering
%         \includegraphics[width=0.6\linewidth]{FIC.png} 
%     \end{minipage}

\end{frame}
\section{Cabeceras}

\begin{frame}[fragile]
\frametitle{Cabeceras}
En la cabecera podemos personalizar el lado izquierdo y el derecho por separado con el paquete \verb|fancyhdr|. Por ejemplo, poner un logo en el izquierdo y un texto en el derecho.
\begin{itemize}
    \item Lado izquierdo $\rightarrow$ \verb|\lhead{}|
    \item Lado derecho $\rightarrow$ \verb|\rhead{}|
\end{itemize}
\end{frame}
\input{content/bibliografía}
\section{Extra}
% Paquetes útiles
\begin{frame}[fragile]
\frametitle{Extra: paquetes útiles}
% placeins
% \FloatBarrier

% xcolor

\end{frame}

% Comandos varios útiles

\begin{frame}[fragile]
\frametitle{Extra: comandos útiles adicionales}

Ocultar el recuadro de los enlaces que aparecen al usar el paquete \verb|hyperref|:
\begin{itemize}
    \item \verb|\hypersetup{hidelinks}|
\end{itemize}

\vspace{0.5cm}

Definir nuestros propios colores:
\begin{itemize}
    \item \verb|\definecolor{micolor}{HTML}{ffdf00}|
    \item Color FIC $\rightarrow$ \verb|\definecolor{FIC}{RGB}{50, 110, 118}|
\end{itemize}

\vspace{0.5cm} 

Cambiar el color de los títulos de las secciones y subsecciones:
\begin{itemize}
    \item \verb|\sectionfont{\color{micolor}}|
    \item \verb|\subsectionfont{\color{micolor}}|
\end{itemize}

\end{frame}

\begin{frame}[fragile]
\frametitle{Extra: comandos útiles adicionales}

Numerar páginas $\rightarrow$ \verb|\pagenumbering{}|
\begin{itemize}
    \item \verb|\pagenumbering{arabic}| $\rightarrow$ 1, 2, 3...
    \item \verb|\pagenumbering{roman}| $\rightarrow$ i, ii, iii...
    \item \verb|\pagenumbering{Roman}| $\rightarrow$ I, II, III...
    \item \verb|\pagenumbering{alph}| $\rightarrow$ a, b, c...
    \item \verb|\pagenumbering{Alph}| $\rightarrow$ A, B, C...
\end{itemize}

\vspace{0.5cm}

No numerar la portada ni el índice:
\begin{itemize}
    \item Eliminar numeración en la portada $\rightarrow$ \verb|\thispagestyle{empty}|
    \item Eliminar numeración en cualquier página antes de comenzar la numeración (índice) $\rightarrow$ \verb|\pagenumbering{gobble}|
\end{itemize}

\end{frame}

\begin{frame}[fragile]
\frametitle{Extra: comandos útiles adicionales}
No numerar las secciones y subsecciones:
\begin{itemize}
    \item No se numera ninguna unidad estructural (ni capítulos ni secciones) 
    
    \hspace{0.2cm} \verb|\setcounter{secnumdepth}{0}|
    
    \item Numerar solo capítulos o secciones 
    
    \hspace{0.2cm} \verb|\setcounter{secnumdepth}{1}|
\end{itemize}

\end{frame}


\end{document}